%% Generated and Refined on 18 March 2026
\documentclass[11pt]{article}

\usepackage[margin=1in]{geometry}
\usepackage{amsmath,amssymb}
\usepackage[hidelinks]{hyperref}
\usepackage[T1]{fontenc}
\usepackage[utf8]{inputenc}

\title{AutoRL: Automated Training of Reinforcement Learning Agents for Predictive Maintenance \\ \Large Version 2.0}
\author{}
\date{18 March 2026}

\begin{document}
\maketitle

\begin{abstract}
This report describes the AutoRL framework: an automated training and evaluation pipeline for reinforcement learning (RL) agents applied to tool wear prediction and maintenance scheduling. We detail four core algorithms (REINFORCE, A2C, PPO, and DQN) and their integration with attention mechanisms. The framework employs grid-based hyperparameter search, checkpoint-based recovery, and multi-run evaluation with statistical aggregation. This document targets technical audiences with basic RL familiarity, emphasizing intuition alongside mathematical formulation.
\end{abstract}

\section{Introduction and Motivation}
Predictive maintenance (PdM) is the problem of determining when a tool should be replaced to minimize downtime while maximizing tool utilization. Manually setting replacement thresholds is suboptimal; RL agents can learn context-aware replacement policies by observing tool wear dynamics and operational constraints.

AutoRL automates the comparison of multiple RL algorithms and attention variants across datasets (SIT, IEEE) to identify the best-performing configuration for PdM tasks. The framework trains agents independently, evaluates them against held-out test data, and produces statistical comparisons with confidence intervals.

\section{Problem Formulation}

\subsection{Markov Decision Process (MDP)}
We model tool wear prediction as a finite-horizon MDP defined by the tuple $\langle \mathcal{S}, \mathcal{A}, P, R, \gamma \rangle$:

\begin{itemize}
\item $\mathcal{S}$: State space (tool wear, operational parameters, environmental factors). \textit{Note: We operate under the Markovian assumption that the current state encapsulates sufficient historical degradation information to predict future wear.}
\item $\mathcal{A} = \{0, 1\}$: Action space ($0 = \text{replace tool}$, $1 = \text{continue operation}$). \textit{Note: This binary formulation effectively frames the predictive maintenance task as an optimal stopping problem.}
\item $P(s'\mid s, a)$: State transition probability.
\item $R(s, a)$: Reward function (tool usage benefits, penalty for violations).
\item $\gamma$: Discount factor for future rewards (detailed in Section 2.2).
\end{itemize}

The agent observes state $s_t$ at each timestep $t$, selects action $a_t \sim \pi(a\mid s_t; \theta)$ where $\pi$ is parameterized by $\theta$, and receives reward $r_t$ and next state $s_{t+1}$.

\subsection{Finite vs. Infinite Horizon and the Discount Factor $\gamma$}
The bounds of the discount factor $\gamma$ depend strictly on the horizon of the MDP:

\textbf{Infinite Horizon ($\gamma \in [0, 1)$):} In non-terminating environments ($T = \infty$), the discount factor must be strictly less than 1. If $\gamma = 1$, the expected cumulative reward $\sum_{t=0}^{\infty} r_t$ may diverge to infinity, destabilizing the value function. 

\textbf{Finite Horizon ($\gamma \in [0, 1]$):} Predictive maintenance is inherently a finite-horizon task; an episode naturally terminates when a tool breaks or is successfully replaced ($T < \infty$). Because the maximum trajectory length is bounded, the cumulative sum $\sum_{t=0}^{T} \gamma^t r_t$ is guaranteed to converge even if $\gamma = 1$. Consequently, setting $\gamma = 1$ is mathematically sound in our formulation, implying the agent values long-term tool survival exactly as much as immediate operation.

\subsection{Objective Function}
The RL objective is to maximize the expected cumulative discounted reward:

\begin{equation}
J(\theta) = \mathbb{E}_{\tau \sim \pi(\theta)} \left[ \sum_{t=0}^{T} \gamma^t r_t \right]
\end{equation}

where $\tau = (s_0, a_0, r_0, \dots, s_T)$ is a trajectory and $T$ is the episode horizon. Different algorithms use different estimators of $\nabla_\theta J(\theta)$.

\section{AutoRL Framework Architecture}

\subsection{Training Pipeline}
AutoRL follows a three-phase pipeline:
\begin{enumerate}
\item \textbf{Training Phase:} For each algorithm--dataset combination, train the agent with gridded hyperparameters (LR, $\gamma$, Attention Type).
\item \textbf{Evaluation Phase:} Load trained models and evaluate on test data (same schema, different files).
\item \textbf{Analysis Phase:} Aggregate results, compute statistics, and generate visualizations.
\end{enumerate}

Checkpoint-based recovery allows interruption and resumption without retraining.

\subsection{Hyperparameter Grid}
The framework supports grid search over:
\begin{itemize}
\item \textbf{Learning Rates:} $\text{LR} \in \{0.001, 0.0001, 0.00001, \dots\}$
\item \textbf{Discount Factors:} $\gamma \in \{0.95, 0.99, 1.0\}$
\item \textbf{Attention Mechanisms:} None, NW (Nadaraya--Watson), TP (Temporal), MH (Multi-Head), SA (Self-Attention)
\end{itemize}

For $N_{\text{algo}} \times N_{\text{LR}} \times N_\gamma \times N_{\text{att}}$ configurations per dataset, the framework trains all combinations automatically.

\section{Reinforcement Learning Algorithms}

\subsection{REINFORCE: Policy Gradient Optimization}
REINFORCE is a foundational policy gradient method. It directly optimizes $J(\theta)$ by sampling full trajectories and computing the gradient:

\begin{equation}
\nabla_\theta J(\theta) = \mathbb{E}_\tau \left[ \sum_{t=0}^{T} \nabla_\theta \log \pi_\theta(a_t\mid s_t) R_t \right]
\end{equation}

where $R_t = \sum_{k=t}^{T} \gamma^{k-t} r_k$ is the return from timestep $t$ onward. The update rule is:

\begin{equation}
\theta \leftarrow \theta + \alpha \nabla_\theta J(\theta)
\end{equation}

\textbf{Advantages:} Unbiased gradient estimates, simple to implement.
\textbf{Disadvantages:} High variance; requires many samples per update. Often benefits from variance reduction techniques.

\subsection{Actor--Critic (A2C): Reducing Variance}
A2C learns two networks: a policy network (actor) and a value network (critic). The critic estimates the baseline value $V(s_t)$ to reduce gradient variance. \textit{Note: Subtracting a state-dependent baseline does not change the expected value of the policy gradient, but it significantly reduces variance, yielding much more stable training.}

\begin{equation}
\nabla_\theta J(\theta) \approx \mathbb{E} \left[ \nabla_\theta \log \pi_\theta(a_t\mid s_t) \cdot (R_t - V_\phi(s_t)) \right]
\end{equation}

The term $(R_t - V_\phi(s_t))$ is the advantage function $A_t$, which measures how much better action $a_t$ is relative to the policy's baseline expectation. The critic minimizes:

\begin{equation}
\mathcal{L}_c(\phi) = \mathbb{E} \left[ (R_t - V_\phi(s_t))^2 \right]
\end{equation}

\textbf{Advantages:} Lower variance than REINFORCE; sample-efficient.
\textbf{Disadvantages:} Two networks require careful tuning; biased if the critic is poorly trained.

\subsection{Proximal Policy Optimization (PPO)}
PPO improves upon A2C by constraining policy updates to a trust region, preventing large detrimental changes:

\begin{equation}
L^{\mathrm{CLIP}}(\theta) = \mathbb{E} \left[ \min\bigl( r_t(\theta) \hat{A}_t, \text{clip}(r_t(\theta), 1-\epsilon, 1+\epsilon) \hat{A}_t \bigr) \right]
\end{equation}

where $r_t(\theta) = \frac{\pi_\theta(a_t\mid s_t)}{\pi_{\theta_{\mathrm{old}}}(a_t\mid s_t)}$ is the probability ratio and $\epsilon$ (typically 0.2) controls the clipping range. \textit{Note: The clipping mechanism explicitly removes the incentive to move too far from the old policy, preventing the "catastrophic policy collapse" often seen in standard gradient methods when step sizes are too large.}

The complete objective function (which is maximized via gradient ascent) combines the clipped surrogate objective, a value function error penalty, and an entropy bonus:

\begin{equation}
J(\theta, \phi) = L^{\mathrm{CLIP}}(\theta) - c_1 L_{VF}(\phi) + c_2 H(\pi_\theta)
\end{equation}

where $H(\pi_\theta)$ is entropy regularization (encourages exploration), and $c_1, c_2$ are weights.
\textbf{Advantages:} Simple, robust, highly stable across varied hyperparameter ranges.
\textbf{Disadvantages:} More tuning knobs than REINFORCE.

\subsection{Deep Q-Network (DQN)}
DQN learns the action-value function $Q(s,a)$ representing expected return for taking action $a$ in state $s$:

\begin{equation}
Q_\pi(s,a) = \mathbb{E} \left[ R_t \mid s_t=s, a_t=a \right]
\end{equation}

The network is trained via temporal difference (TD) learning:

\begin{equation}
\mathcal{L}(\theta) = \mathbb{E} \left[ \left( r + \gamma \max_{a'} Q_{\theta_{\mathrm{target}}}(s', a') - Q_\theta(s, a) \right)^2 \right]
\end{equation}

Experience replay and a separate target network stabilize training. \textit{Note: Standard DQN suffers from overestimation bias due to the $\max$ operator in the target calculation, which is theoretically mitigated in advanced implementations via Double Q-Learning (DDQN).}
\textbf{Advantages:} Naturally handles discrete actions; can learn off-policy from historical data.
\textbf{Disadvantages:} Q-function overestimates; restricted to discrete action spaces.

\section{Attention Mechanisms for RL}
Attention mechanisms allow agents to focus on relevant parts of their observation. In PdM, this means dynamically weighting sensors, time steps, or features across the tool's lifespan.

\subsection{General Attention Concept}
Given a set of values $V = \{v_1, \ldots, v_n\}$ and queries $q$, attention computes a weighted sum:

\begin{equation}
\mathrm{Attention}(q, V) = \sum_{i=1}^{n} \alpha_i v_i
\end{equation}

where attention weights $\alpha_i$ sum to 1 and are computed via:

\begin{equation}
\alpha_i = \frac{\exp(s(q, v_i))}{\sum_{j=1}^{n} \exp(s(q, v_j))}
\end{equation}

Here $s(q, v_i)$ is a scoring function (e.g., dot product, learned similarity).

\subsection{Nadaraya--Watson Attention}
Nadaraya--Watson (NW) is a kernel-based attention using Gaussian kernels. For temporal features $x_t$ and current query $q$:

\begin{equation}
\begin{aligned}
\alpha_i^{(t)} &= \frac{K(x_i, q_t)}{\sum_j K(x_j, q_t)},\\
K(x, q) &= \exp\left( - \frac{\lVert x - q \rVert^2}{2\sigma^2} \right)
\end{aligned}
\end{equation}

This creates smooth, non-parametric weighting.
\textbf{Useful for:} Detecting similarity between past and current observations based on Euclidean distance.

\subsection{Temporal Attention}
Temporal attention weights past timesteps; effective for sequences:

\begin{equation}
\alpha_i^{(t)} = \frac{\exp\left( \frac{W_q q_t \cdot (W_k x_i)^T}{\sqrt{d}} \right)}{\sum_j \exp\left( \frac{W_q q_t \cdot (W_k x_j)^T}{\sqrt{d}} \right)}
\end{equation}

where $W_q, W_k$ are learnable projections and $d$ is dimension. \textit{Note: The scaling by $\sqrt{d}$ stabilizes gradients by preventing the dot products from growing excessively large, which would otherwise push the softmax function into saturated regions and cause vanishing gradients.}
\textbf{Useful for:} Long-range dependencies in wear trajectories.

\subsection{Multi-Head Attention}
Multi-head attention (MH) runs multiple attention heads in parallel:

\begin{equation}
\mathrm{MultiHead}(q, V) = \mathrm{Concat}(\mathrm{head}_1, \ldots, \mathrm{head}_h)W^O
\end{equation}

where each $\mathrm{head}_i = \mathrm{Attention}(q, V)_i$ attends to different representational subspaces. This captures multiple types of degradation relationships simultaneously.

\subsection{Self-Attention}
Self-attention (SA) weights observations relative to each other within a sequence:

\begin{equation}
\mathrm{SelfAttn}(X) = \mathrm{softmax}\left( \frac{X W_q (X W_k)^T}{\sqrt{d}} \right) (X W_v)
\end{equation}

where $X$ is the sequence of observations and $W_q, W_k, W_v$ are weight matrices. Self-attention enables the agent to relate tool wear at time $t$ to wear at time $j$ directly.
\textbf{Advantage:} Captures long-range temporal dependencies without explicit temporal decay mechanisms.
\textbf{Disadvantage:} Computational complexity scales quadratically $\mathcal{O}(T^2)$ with respect to sequence length $T$.

\section{Evaluation Methodology}

\subsection{Evaluation Metrics}
After training, agents are evaluated on test data from the same schema. For each trajectory, three metrics are computed:
\begin{itemize}
\item \textbf{Tool Usage:} Percentage of timesteps where action is ``continue'' (higher is better, max 100\%)
\item \textbf{Lambda:} Mean deviation of tool wear at replacement from target region (lower is better)
\item \textbf{Threshold Violations:} Count of timesteps where tool wear exceeds safety threshold (lower is better)
\end{itemize}

\subsection{Composite Evaluation Score}
A weighted evaluation score combines these metrics:

\begin{equation}
\mathrm{Score} = 0.50 \cdot T_{\mathrm{usage}} + 0.30 \cdot (1 - \mathrm{norm}_{\lambda}) + 0.20 \cdot V_{\mathrm{violations}}
\end{equation}

where $T_{\mathrm{usage}}$ is tool usage normalized to $[0, 1]$, $\mathrm{norm}_{\lambda}$ is normalized deviation, and $V_{\mathrm{violations}}$ is the violation count normalized via $\frac{1}{1+v}$. \textit{Note: This inverse normalization elegantly converts a "lower-is-better" penalty into a "higher-is-better" reward signal, ensuring the agent can consistently maximize the overall scalar score.}

\subsection{Statistical Analysis}
For fair comparison, we compute aggregated statistics:

\begin{equation}
\begin{aligned}
\mu &= \frac{1}{n} \sum_{i} \mathrm{Score}_i, \\
\sigma &= \sqrt{\frac{1}{n-1} \sum_{i} (\mathrm{Score}_i - \mu)^2}
\end{aligned}
\end{equation}

\begin{equation}
\begin{aligned}
\mathrm{SEM} &= \frac{\sigma}{\sqrt{n}}, \\
\text{95\% CI} &= [\mu - 1.96 \cdot \mathrm{SEM},\; \mu + 1.96 \cdot \mathrm{SEM}]
\end{aligned}
\end{equation}

Confidence intervals account for sample variability; overlapping CIs indicate no statistically significant difference between configurations.

\section{Experimental Design}
The AutoRL pipeline executes:
\begin{enumerate}
\item Train all algorithm--attention hyperparameter combinations on each training file
\item Evaluate trained models on all test files (same schema)
\item Compute per-model and per-algorithm statistics
\item Generate heatmaps and multi-panel analysis reports
\end{enumerate}

This enables the empirical identification of:
\begin{itemize}
\item \emph{Best overall algorithm:} Highest average score across all configurations
\item \emph{Best attention type:} Largest improvement with attention enabled compared to baseline
\item \emph{Most stable algorithm:} Lowest variance across runs and random seeds
\item \emph{Dataset-specific insights:} Which algorithms or attention heads excel on which specific datasets
\end{itemize}

\end{document}